\subsection{Objetivos}

\subsubsection{Objetivo General}
Analizar y caracterizar de manera integral el Proceso de Matrícula Regular de Estudiantes de la UNI, abarcando la articulación de sus cuatro subprocesos principales para identificar su cadena de valor, diferenciar la dimensión procedimental de la tecnológica y diagnosticar los puntos de dolor en su estado actual (As-Is).

\subsubsection{Objetivos Específicos}
\begin{itemize}
\item Describir el esquema de transferencia de entradas y salidas (handoffs) y los bucles de control (como los bucles de excepción por solicitudes o de invalidación por sanción) entre los subprocesos.
\item Caracterizar detalladamente los insumos, actividades, actores, recursos, reglas de negocio y salidas de los subprocesos de Acreditación, Selección, Verificación y Confirmación.
\item Diferenciar las responsabilidades normativas (reglas institucionales y de los actores) de las capacidades del sistema de información (motor tecnológico DIRCE/SIGA-ORCE).
\item Levantar y diagnosticar los problemas frecuentes en cada etapa, tales como la desincronización de pagos, saturación del portal de matrícula, cruces de horarios y las demoras operativas en el cierre de actas.
\end{itemize}